Each of $x$ and $x+1$ must have at least two distinct prime factors. Indeed, a prime power dividing $y(y+1)$ divides one of the coprime factors $y,y+1$, contrary to the hypothesis. The smallest consecutive positive integers with this property are $14,15$, so $x\ge14$. Also $y>x$, since $y(y+1)$ is a positive multiple of $x(x+1)$ and equality $y=x$ is forbidden.

For $x=14$, testing $y=15,16,17,18,19$ gives no multiple of $210$, while $y=20$ gives $20\cdot21=420$. Neither $14$ nor $15$ divides $20$ or $21$, so this works with $x^2+y^2=196+400=596$.

The values $x=15,16,17$ fail the prime-factor condition because $16$ is a prime power and $17$ is prime. For $x\ge18$, $x^2+y^2\ge18^2+19^2=685>596$. Therefore the minimum is $\boxed{596}$.
